\chapter{Theoretische Grundlagen}
\label{kap:theorie}

\section{Convolutional Neural Networks}
\label{sec:cnn}

\subsection{Grundprinzipien}
\label{sec:prinzipien}

Convolutional Neural Networks beruhen auf drei Prinzipien: \emph{lokale
rezeptive Felder}, \emph{geteilte Gewichte} und \emph{räumliche
Unterabtastung}. Sie sind keine Implementierungsdetails, sondern kodieren
strukturelles Vorwissen über Bilddaten in die Netzarchitektur -- Lokalität,
Stationarität und Toleranz gegenüber kleinen Verschiebungen
\parencite{lecun1998,goodfellow2016}.

\absatz{Lokale rezeptive Felder und geteilte Gewichte.} Statt jedes
Ausgabeneuron mit allen $H \cdot W$ Eingangspixeln zu verbinden, beschränkt eine
Faltungsschicht die Summation auf eine $k \times k$-Nachbarschaft und verwendet
denselben Kernel an allen Positionen. Die Parameterzahl einer Schicht wird
dadurch unabhängig von der Bildgröße, und dasselbe Muster wird überall erkannt,
auch wenn es in den Trainingsdaten nur an wenigen Positionen vorkommt. Der
belastbare Vorteil liegt in dieser \emph{Stichprobeneffizienz}, nicht in der
Kapazitätsbeschränkung: \textcite{zhang2017rethinking} zeigen, dass gängige
Faltungsnetze zufällige Beschriftungen vollständig memorieren können und dennoch
generalisieren. Die Operation ist \emph{translationsäquivariant} -- eine
verschobene Eingabe erzeugt eine identisch verschobene Ausgabe --, was durch die
näherungsweise Stationarität natürlicher Bildstatistik gerechtfertigt ist. Die Faltung
 ist dabei translationsäquivariant, während nachfolgendes Downsampling die räumliche
 Präzision reduziert und dadurch eine gewisse Robustheit gegenüber kleinen Verschiebungen 
 erzeugen kann

\absatz{Räumliches Subsampling -- die für diese Arbeit entscheidende Stufe.} Für
die Objekterkennung ist Invarianz gegenüber kleinen Verschiebungen erwünscht.
Sie entsteht durch stufenweise Reduktion der räumlichen Auflösung -- klassisch
durch Max-Pooling, in modernen Architekturen wie YOLO26 überwiegend durch
Faltungen mit Schrittweite $s = 2$. Erhalten bleibt dabei nur, \emph{ob} ein
Merkmal im Fenster auftrat, nicht seine exakte Position. Nach $L$ solchen Stufen
beträgt die \textbf{Gesamtschrittweite} $S = \prod_\ell s_\ell = 2^L$.

Der Preis ist ein Verlust an Ortsinformation, und der wird bei kleinen Objekten
kritisch. Eine Zelle der Merkmalsebene mit Gesamtschrittweite $S$ entspricht
einem Bildbereich von $S \times S$ Pixeln; ein Objekt mit Durchmesser $d$ belegt
dort nur $d/S$ Zellen. Für \(d<S\) erstreckt sich das Objekt im Mittel über weniger
 als eine Gitterzelle und ist damit nur durch sehr wenige räumliche Positionen repräsentiert.
  Eine zuverlässige Lokalisierung wird dadurch zunehmend erschwert. Für die feinste
Standard-Detektionsebene P3 in YOLO26 gilt $S = 8$ -- das kleinste Rohrprofil in
\code{test} misst 6,4 Pixel und unterschreitet diese Grenze. „Das Verhältnis \(d/S\) 
dient in dieser Arbeit als zentrale geometrische Kenngröße für die räumliche Repräsentierbarkeit
 kleiner Objekte es begründet
die Multi-Skalen-Fusion moderner Detektoren (Abschnitt~\ref{sec:neck}) und die
drei möglichen Gegenmaßnahmen, die alle dasselbe Verhältnis vergrößern: Eine
höhere Eingabeauflösung und Kachelung erhöhen $d$, ein zusätzlicher
Detektionskopf senkt $S$ (Abschnitt~\ref{sec:p2kopf}).

\subsection{Merkmalshierarchie und Objektdetektion}
\label{sec:hierarchie}

Auf jede Faltung folgt eine nichtlineare Aktivierung -- heute überwiegend die
Rectified Linear Unit $f(x) = \max(0,x)$ oder Varianten wie SiLU
\parencite{krizhevsky2012} --, ohne die das Netz nur lineare Abbildungen lernen
könnte. Aus der Stapelung entsteht ein hierarchischer Aufbau der
Repräsentationen: Frühe Schichten reagieren auf Kanten und Farbverläufe, mittlere
auf Texturen und Teilstrukturen, tiefe auf semantische Objektbestandteile
\parencite{zeiler2014}. Diese datengetriebene Merkmalsextraktion tritt an die
Stelle des manuell entworfenen Feature Engineering und gilt als wesentlicher
Grund für die Leistungssprünge tiefer Netze \parencite{lecun2015}. Am Ende
klassischer Architekturen fassen vollvernetzte Schichten die Merkmale zusammen;
in modernen Detektoren sind sie vollständig durch Faltungen ersetzt
(Abschnitt~\ref{sec:neck}). Residualverbindungen \parencite{he2016} ermöglichen
das stabile Training sehr tiefer Netze und sind Standardbaustein nahezu aller
CNN-Backbones.

Für diese Arbeit ist nicht die Klassifikation relevant, sondern die
\textbf{Objektdetektion}: das gleichzeitige Lokalisieren und Klassifizieren
mehrerer Objekte. Zweistufige Detektoren erzeugen zunächst Regionsvorschläge und
klassifizieren diese anschließend \parencite{ren2015}; einstufige Detektoren
formulieren die Detektion als einmaliges Regressionsproblem und sagen in einem
Durchlauf direkt Boxen und Klassenwahrscheinlichkeiten vorher
\parencite{redmon2016,liu2016ssd}, was Echtzeitfähigkeit bei geringen
Genauigkeitseinbußen ermöglicht. Klassisch nachgeschaltet ist die
Non-Maximum Suppression (NMS), die redundante überlappende Detektionen anhand
ihrer Intersection over Union unterdrückt -- bei dicht gepackten gleichartigen
Objekten eine potenzielle Fehlerquelle.

\section{YOLO26}
\label{sec:yolo26}

\subsection{Einordnung}

\enquote{You Only Look Once} wurde 2016 als erster echtzeitfähiger einstufiger
Detektor vorgestellt \parencite{redmon2016}; ab YOLOv8 übernahm Ultralytics die
Weiterentwicklung und etablierte ankerfreie, entkoppelte Detektionsköpfe
\parencite{sapkota2025}. YOLO26 wurde im September 2025 vorgestellt und
unterscheidet sich von YOLO11 durch vier koordinierte Änderungen: ein
Dual-Head-Design für NMS-freie Inferenz, den Wegfall der Distribution Focal Loss,
ProgLoss zusammen mit STAL sowie den MuSGD-Optimierer
\parencite{jocher2026yolo26,sapkota2025}.

Für diese Arbeit ist davon eine Änderung einschlägig: Die
\emph{Small-Target-Aware Label Assignment} garantiert positive Labelzuweisungen
auch für sehr kleine Objekte -- genau das Problem klassischer
Zuweisungsstrategien, das Abschnitt~\ref{sec:ioukleinobjekt} quantifiziert.

\subsection{Backbone, Neck, Head und die Merkmalspyramide}
\label{sec:neck}

YOLO26 folgt der Dreiteilung \textbf{Backbone $\to$ Neck $\to$ Head}. Das
Backbone reduziert die räumliche Auflösung stufenweise per Stride-2-Faltungen
($1024 \to 512 \to 256 \to 128 \to 64 \to 32$ bei $1024$\px{} Eingabe), während
die Kanalzahl wächst -- räumliche Präzision wird gegen semantische Tiefe
getauscht.

\absatz{Neck.} Das Neck verbindet die Auflösungsstufen in beide Richtungen
(FPN/PAN-Prinzip, \cite{lin2017fpn}): Tiefe, semantisch reiche Merkmalskarten
werden hochskaliert und mit den feinen, ortsgenauen Karten früherer Stufen
zusammengeführt, und umgekehrt. Der Head erhält dadurch mehrere Ausgabeebenen,
die nach Objektgröße zuständig sind. Die Standardkonfiguration liefert drei
Ebenen P3/P4/P5 mit den Gesamtschrittweiten 8/16/32; zusätzlich existiert eine
Konfiguration mit einer vierten, feineren Ebene \textbf{P2 mit Schrittweite~4}
(Abbildung~\ref{abb:pyramide}). Diese Variante wird hier durchgängig verwendet;
Abschnitt~\ref{sec:p2kopf} begründet, warum.

\begin{figure}[htbp]
\centering
\begin{tikzpicture}[x=1cm,y=1cm]
  % ---- Merkmalsebenen als Quadrate, Seitenlaenge proportional zur Gittergroesse
  \draw[fill=red!8,dashed,thick,red!60!black] (0,0) rectangle (4,4);
  \draw[step=0.5,gray!35,very thin] (0,0) grid (4,4);
  \draw[fill=blue!8] (4.7,0) rectangle (6.7,2);
  \draw[step=0.5,gray!35,very thin] (4.7,0) grid (6.7,2);
  \draw[fill=blue!8] (7.4,0) rectangle (8.4,1);
  \draw[step=0.5,gray!35,very thin] (7.4,0) grid (8.4,1);
  \draw[fill=blue!8] (9.1,0) rectangle (9.6,0.5);

  \node[below,align=center,font=\scriptsize] at (2,-0.2)
       {\textbf{P2} $\;$ Stride 4\\$256\times256$\\65\,536 Positionen};
  \node[below,align=center,font=\scriptsize] at (5.7,-0.2)
       {\textbf{P3} $\;$ Stride 8\\$128\times128$\\16\,384};
  \node[below,align=center,font=\scriptsize] at (7.9,-0.2)
       {\textbf{P4} $\;$ 16\\$64^2$\\4\,096};
  \node[below,align=center,font=\scriptsize] at (9.35,-0.2)
       {\textbf{P5} $\;$ 32\\$32^2$\\1\,024};

  \node[above,font=\scriptsize,red!60!black,align=center] at (2,4.15)
       {\emph{zusätzlich in} \code{-p2}, \emph{hier verwendet}};
  \node[above,font=\scriptsize,align=center] at (7.0,4.15)
       {Standardkonfiguration P3--P5};
  \draw[decorate,decoration={brace,amplitude=4pt},gray!60!black]
       (4.7,4.05) -- (9.6,4.05);
  \node[font=\scriptsize,align=left,anchor=west] at (10.1,2.6)
       {kleine\\Objekte};
  \node[font=\scriptsize,align=left,anchor=west] at (10.1,1.2)
       {große\\Objekte};
  \draw[-{Latex[length=1.6mm]},gray!60!black] (10.35,2.3) -- (10.35,1.6);
\end{tikzpicture}
\caption[Merkmalsebenen des Detektionskopfes]{Merkmalsebenen des
Detektionskopfes bei einer Eingabe von $1024\times1024$ Pixeln. Die
Seitenlänge der Quadrate ist proportional zur Gittergröße. Jede Zelle einer
Ebene mit Schrittweite $S$ entspricht einem Bildbereich von $S \times S$ Pixeln
und sagt Boxkoordinaten sowie Klassen-Scores vorher. Die Standardkonfiguration
nutzt \textbf{P3--P5}; die hier trainierten Modelle verwenden zusätzlich die
links rot gestrichelt eingezeichnete Ebene \textbf{P2} mit Schrittweite~4. Die
Abbildung zeigt zugleich, warum diese eine Ebene so viel ausmacht: P2 allein
enthält mehr Prädiktionspositionen als P3, P4 und P5 zusammen. Genau diese
Abhängigkeit von der Schrittweite macht Abschnitt~\ref{sec:ioukleinobjekt} zum
Kern des Auflösungsproblems.}
\label{abb:pyramide}
\end{figure}

\absatz{Head.} Der Kopf ist vollständig faltungsbasiert: kleine
Faltungsstapel, die in $1\times1$-Faltungen münden. Eine $1\times1$-Faltung ist
äquivalent zu einer vollvernetzten Schicht über die Kanaldimension, wird aber
\textbf{positionsgeteilt} angewendet -- dieselbe Operation verarbeitet die Kanäle
jeder Gitterzelle separat. Das erhält die Translationsäquivarianz bis in die
Ausgabeschicht und macht das Netz unabhängig von einer festen Eingabegröße. Der
Kopf ist zusätzlich
\emph{entkoppelt} (getrennte Zweige für Klassifikation und Boxregression) und
\emph{ankerfrei}: Jede Gitterzelle regressiert direkt die Abstände zu den vier
Boxrändern. Damit ist die Zahl der Gitterzellen -- und nicht die Zahl der
Ankerboxen -- die Größe, die bestimmt, wie fein Objekte überhaupt aufgelöst
werden können.

\subsection{Der P2-Detektionskopf}
\label{sec:p2kopf}

Alle in dieser Arbeit trainierten Modelle sind \code{yolo26s-p2}. Diese Wahl ist
kein Gegenstand der Experimente, sondern deren Konstante; sie ist an dieser
Stelle zu begründen, weil sie unmittelbar an der Auflösungsgrenze hängt, die
Abschnitt~\ref{sec:ioukleinobjekt} herleitet.

\absatz{} Der Unterschied liegt allein im Neck und im Kopf. Die
Standardkonfiguration führt den Aufwärtspfad des Feature-Pyramid-Netzwerks
\parencite{lin2017fpn} bis P3 und schickt von dort einen Abwärtspfad über P4 zu
P5; die \code{-p2}-Variante setzt den Aufwärtspfad um eine Stufe fort, verknüpft
ihn mit der P2-Karte des Backbones und startet den Abwärtspfad eine Ebene tiefer.
Der Detektionskopf erhält vier statt drei Eingänge. Weil die Eingangsbreite des
Kopfes seine internen Kanalzahlen bestimmt und der neue, feinste Eingang schmaler
ist als der bisherige, hat die Variante sogar \emph{weniger} Parameter als das
Basismodell: 9\,663\,464 gegen 9\,948\,638, also $-2{,}9\,\%$. Alle Kennzahlen
dieses Absatzes stammen aus der Modellzusammenfassung, die die Werkzeugkette beim
Aufbau der beiden Netze ausgibt; Speicher und Laufzeit sind an den eigenen
Trainingsläufen gemessen. Sie ist kein
größeres, sondern ein \emph{feiner aufgelöstes} Netz.

\absatz{}Die Prädiktionspositionen sind die Summe der
Gitterzellen aller Detektionsebenen und damit aus der Eingabegröße abzählbar: Bei
$1024\times1024$ vervierfachen sie sich von $128^2+64^2+32^2 = 21\,504$ auf
$256^2 + 21\,504 = 87\,040$. Der Rechenaufwand steigt dagegen nur um 17\,\%
(22,7 auf 26,6 GFLOPs bei $640^2$, laut Modellzusammenfassung), weil die zusätzlichen Blöcke
auf der feinsten Ebene schmal sind. Teuer ist der \emph{Speicher}: Er wird von
den Aktivierungen dominiert, und eine Merkmalskarte auf P2 hat viermal so viele
Positionen wie eine auf P3. Gemessen belegt das Training bei \code{imgsz}=1024
mit Batch~8 rund 12\,GB statt etwa 6\,GB. Das ist der eigentliche Preis der
Variante und der Grund, warum sie nicht der Standard ist.

\absatz{Nutzen} Der Nutzen lässt sich vor jedem Training abschätzen.
Nach Gleichung~\eqref{eq:zellen} beträgt die erwartete Zahl der Gitterzellen,
deren Zentrum in ein kreisförmiges Objekt vom Durchmesser $d$ fällt und die damit
überhaupt als positive Zuweisung in Frage kommen, $\mathbb{E}[N] = \pi d^2/(4S^2)$;
der Übergang von $S = 8$ auf $S = 4$ vervierfacht diesen Wert.
Tabelle~\ref{tab:p2zellen} setzt ihn für die tatsächlich auftretenden
Objektgrößen ein. Entscheidend ist nicht der Mittelwert, sondern die Spalte
\enquote{min}: Bei $S = 8$ hat das kleinste Objekt in \code{test} im
ungünstigsten Fall \emph{genau eine} Kandidatenzelle -- eine einzige Position,
die es finden kann --, bei $S = 4$ sind es vier. Die Zuweisungsgrenze
$d < S\sqrt{2}$ sinkt von 11,3\px{} auf 5,7\px{} und liegt damit vollständig
unterhalb der hier auftretenden Objektgrößen.

\begin{table}[htbp]
\centering
\caption[Kandidatenzellen je Objekt]{Erwartete Zahl der Gitterzellen, deren
Zentrum in das Objekt fällt, nach Gleichung~\eqref{eq:zellen}, bei nativem
Maßstab. Die Spalte \enquote{min} nennt den ungünstigsten Fall über alle Lagen
des Objektzentrums im Gitter; sie hat keine geschlossene Form und ist durch
Absuchen der Lagen auf einem feinen Raster bestimmt.}
\label{tab:p2zellen}
\footnotesize
\begin{tabular}{@{}rlrrrr@{}}
\toprule
$d$ (px) & Bedeutung & \multicolumn{2}{c}{P3 ($S=8$)} &
\multicolumn{2}{c}{P2 ($S=4$)} \\
 & & Mittel & min & Mittel & min \\
\midrule
11,6 & p05 des synthetischen Datensatzes & 1,7 & 1 & 6,6 & 4 \\
12,2 & kleinstes Objekt in \code{test}   & 1,8 & 1 & 7,3 & 4 \\
19,6 & Median \code{test}                & 4,7 & 4 & 18,9 & 16 \\
22,5 & Median des synthetischen Datensatzes & 6,2 & 4 & 24,9 & 21 \\
27,4 & Median \code{val}                 & 9,2 & 8 & 36,9 & 32 \\
\bottomrule
\end{tabular}
\end{table}

\absatz{} Ein Detektor, dem die
Auflösung fehlt, übersieht kleine Objekte unabhängig davon, wie gut seine
Trainingsdaten sind. Jede Verbesserung der Datenerzeugung würde dann teilweise
nur eine Architekturschwäche kompensieren, und der gemessene Effekt wäre nicht
der Effekt der Daten. Die Ablationen der Kapitel~\ref{kap:auswertung} laufen
deshalb auf der feiner aufgelösten Variante -- nicht, weil sie besser
\emph{aussieht}, sondern weil sie die Störgröße entfernt, die den Datengewinn
sonst überlagert.

\section{Anwendung auf gestapelte Bauzaunprofile}
\label{sec:anwendung}

Ein YOLO-Modell wird per Transfer Learning von COCO-vortrainierten Gewichten
\parencite{yosinski2014} auf eine einzige Objektklasse feinabgestimmt. Zwei
Eigenschaften von YOLO26 passen dabei unmittelbar zum Problem: STAL und ProgLoss
verbessern die Erkennung kleiner Objekte, und die NMS-freie Inferenz eliminiert
einen bei dicht gepackten gleichartigen Objekten systematisch fehleranfälligen
Hyperparameter -- ein zu niedriger IoU-Schwellwert unterdrückt korrekte
Detektionen benachbarter Rohre, ein zu hoher lässt Duplikate zu.

Drei Grenzen bleiben und strukturieren den Praxisteil. Die \textbf{Auflösung}
(Abschnitt~\ref{sec:prinzipien}) wird über den P2-Kopf und die Inferenzauflösung
adressiert (Abschnitte~\ref{sec:p2kopf} und~\ref{sec:imgsz}); die
\textbf{Verdeckung} -- verkippte Zäune verdecken benachbarte Profile, und
Gitterknoten, Löcher im Maschendraht und Schattenspalten drohen als
Falschpositive -- ist Gegenstand der Ablation in Abschnitt~\ref{sec:ablationB};
und der \textbf{Datenbedarf} wird nicht durch Handannotation gedeckt, sondern
durch eine Synthese-Pipeline, deren Entwurf die Kapitel~\ref{kap:augtheorie}
bis~\ref{kap:trainingszeit} füllt.
